\section{The shared problem: too many interactions}

\subsection{What does a block mean?}
Suppose $N$ sources $y_j$ influence $N$ targets $x_i$ through
\begin{equation}
 u_i=\sum_{j=1}^{N}K(x_i,y_j)q_j,\qquad i=1,\ldots,N.
 \label{eq:pairwise}
\end{equation}
In matrix notation, $u=Kq$. A direct implementation stores and applies an $N\times N$ table. \textbf{Applying} the operator means computing $Kq$; \textbf{solving} means finding $q$ such that $Kq=b$. The two tasks share compressed representations, but solving also requires safe elimination and pivoting.

A \emph{block} is a rectangular subset $B=K(X,Y)$ of that table, determined by a row index set $X$ and a column index set $Y$. Its \emph{numerical $\eps$-rank} is
\begin{equation}
 r_\eps(B)=\min\{r:\|B-B_r\|_2\leq \eps\|B\|_2\}.
 \label{eq:numericalrank}
\end{equation}
By the Eckart--Young theorem this is the number of singular values above $\eps\sigma_1(B)$. If $B_r=UV^{\T}$ with $U\in\mathbb R^{m\times r}$ and $V\in\mathbb R^{n\times r}$, one stores $r(m+n)$ rather than $mn$ leading numbers. For complex data, use $UV^*$ where appropriate. \textbf{Exact algebraic rank and numerical $\eps$-rank are different}: the DFT example in Section~\ref{sec:butterfly} makes the distinction visible.

A short representation is therefore not just ``a small rank''. It is a particular stored object, at a stated tolerance and norm, on a stated block geometry. Rank is not an invariant of a PDE alone.

\subsection{A complete example: why some blocks are short}
\label{sec:logexample}
Take $K(x,y)=\log(x-y)$, with $x>0$ and $|y|\le a<x$. Then
\begin{equation}
 \log(x-y)=\log x+\log(1-y/x)
 =\log x-\sum_{q=1}^{\infty}\frac{1}{q}(y/x)^q.
 \label{eq:laplaceseries}
\end{equation}
Every term is a product of a function of $x$ and a function of $y$. Truncating after $q=p$ produces at most $p+1$ separable terms. If $\rho=a/x<1$, the discarded tail obeys
\begin{equation}
 \left|\sum_{q=p+1}^{\infty}\frac{1}{q}(y/x)^q\right|
 \leq \sum_{q=p+1}^{\infty}\frac{\rho^q}{q}
 \leq \frac{\rho^{p+1}}{(p+1)(1-\rho)}.
 \label{eq:logbound}
\end{equation}
That is a complete small compression theorem: \textbf{separation limits $\rho$; a power series gives the representation; the tail bounds the error}. Other kernels need other expansions, bounds and assumptions. Smoothness alone is not a universal compression certificate.

\subsection{Why trees enter}
A single box cannot be far from every target. Hierarchical grouping handles this: small boxes describe local interactions exactly, while larger boxes describe distant ones by short summaries. A parent may reuse information from its children. A tree is therefore an \textbf{accounting device} for covering interactions without treating them all at the most expensive scale; it does not itself prove that a block is low rank.

The rest of the report asks the same five questions of every method.

\section{The five design choices: a common grammar, not a universal recipe}
\label{sec:dials}

\begin{center}
\begin{tikzpicture}[
 font=\small,box/.style={draw=teal,fill=light,rounded corners,align=center,minimum height=13mm},
 every node/.style={align=center}]
\node[box,text width=2.9cm] (p) {1. Partition\\[-2pt]\scriptsize which groups?};
\node[box,text width=2.9cm,right=4mm of p] (a) {2. Interaction pattern\\[-2pt]\scriptsize which blocks?};
\node[box,text width=2.9cm,right=4mm of a] (r) {3. Representation\\[-2pt]\scriptsize which summary?};
\node[box,text width=2.9cm,below=5mm of r] (s) {4. Sampling\\[-2pt]\scriptsize how is it learned?};
\node[box,text width=2.9cm,left=4mm of s] (g) {5. Algebra\\[-2pt]\scriptsize apply or solve?};
\draw[-{Stealth[length=2mm]},thick,teal] (p)--(a);
\draw[-{Stealth[length=2mm]},thick,teal] (a)--(r);
\draw[-{Stealth[length=2mm]},thick,teal] (r)--(s);
\draw[-{Stealth[length=2mm]},thick,teal] (s)--(g);
\end{tikzpicture}
\end{center}

The code calls these choices \textbf{dials}, because several can be swapped while preserving an interface. Mathematically, they are \textbf{design choices}. Each answers a different question; none supplies the evidence needed by the next one.

\subsection{Choice 1: partition and interaction geometry}
\textbf{Question.} Which rows and columns belong together, in which coordinate system and in which order? A spatial binary tree, quadtree or octree groups points by position. A butterfly needs two trees, one for positions and one for frequencies. A transform may instead tree an angle and a spectral degree. The DFT additionally needs bit-reversed index order before residue classes become contiguous intervals.

A partition fixes the candidate blocks, the tree depth, the leaf cost and the opportunities to reuse a parent's work. It does \emph{not} prove that those blocks are compressible. The same matrix with natural, one-sided bit-reversed and two-sided bit-reversed order can have very different numerical $\eps$-ranks (the DFT/FFT ordering experiment). Nor does a smaller leaf or root automatically reduce total work: it also changes the number of blocks and transfers.

\textbf{Failure signal.} Ranks grow because the grouping no longer matches the operator's geometry, or a representation that exists mathematically cannot be reached by the current coordinates. The DFT/FFT experiment responds with residue classes and bit reversal; the spherical-harmonic experiment uses trees in colatitude $\theta$ and degree $l$; the Oseen experiment tests anisotropic geometry near a wake.

\subsection{Choice 2: admissibility and the interaction list}
\textbf{Question.} Which block pairs are compressed, and which are kept explicitly? In one-tree \emph{weak} admissibility, every off-diagonal pair at a parent level is compressed, including touching boxes; leaf and diagonal data remain explicit. In \emph{strong} admissibility, only sufficiently separated pairs are compressed and touching neighbours remain in the near field. The precise distance ratio is an implementation choice. The near/far split determines both approximation work and the exact sparse remainder.

Butterfly \emph{complementary} pairing is different: it is not a Euclidean-distance test. At stage $\ell$ it pairs a row node at depth $\ell$ with a column node at depth $L-\ell$. This produces a sequence of tilings whose row and column widths trade off. Directional rules enlarge or anisotropically deform the near list.

Admissibility does \emph{not} establish the approximation error of an admissible pair, the stability of an elimination, or a cheaper whole algorithm. Strong admissibility, for example, can keep so many near blocks that storage and elimination work rise even when its ranks are smaller (the Oseen wake experiment).

\textbf{Failure signal.} A supposedly admissible block is not low rank, or rank improves but the exact near field dominates. The interaction-geometry experiment changes the rule; the high-frequency-arc experiment uses complementary storage only for separated blocks; the screening experiment allows negligible near blocks to move back into compression.

\subsection{Choice 3: representation and compression}
\textbf{Question.} What mathematical object replaces an admissible block? Three families recur:
\begin{itemize}
 \item \textbf{Analytic summaries.} Multipole/local expansions, plane-wave signatures and their truncation orders. Translation acts on these summaries.
 \item \textbf{Algebraic factors.} $B\approx UV^*$, or an interpolative decomposition (ID) that selects skeleton columns $S$ and represents redundant columns $R$ by $B_{:,R}\approx B_{:,S}T$.
 \item \textbf{Nested representations.} Bases or transfers shared between levels, as in $\mathcal H^2$, HBS and butterfly factorization. Nesting changes the entry count, not merely the rank of an isolated block.
\end{itemize}
The representation determines what is stored and what a transfer means. It does \emph{not} specify how to obtain the factors, how to assemble them, or whether elimination will be stable. An SVD describes an accurate block space; an ID additionally selects rows or columns useful in elimination; neither chooses an FMM translation or a safe pivot.

\textbf{Failure signal.} A single whole-block low-rank factor needs rank growing with frequency (the high-frequency-arc experiment), independent tiles need too many separate bases (butterfly), or block-valued and multi-order kernels do not fit a scalar single-order representation (Stokes and Matérn sampling).

\subsection{Choice 4: sampling and construction}
\textbf{Question.} Which data determine the representation, and what entries or operator products must be read? Exact sampling reads true far-block rows. Proxy sampling evaluates auxiliary points on a circle, sphere, annulus, lifted surface or volume. Randomized construction uses matrix--vector products or sketches. Entry-sampled butterfly constructors use selected entries and require their own coherence hypotheses.

Sampling is both a \textbf{mathematical sufficiency hypothesis} and a \textbf{cost model}. A proxy is sufficient only if observations on the auxiliary set determine the fields that matter at the true far points. A cheap factor is not cheap to build if its constructor still reads all $N^2$ entries. Conversely, an exact far-block ID can be accurate yet interact badly with a later complex elimination (complex Helmholtz elimination).

Sampling does \emph{not} follow from kernel smoothness, does not certify a solve, and does not inherit an arbitrary block's relative scale safely. In exponentially small blocks, a per-tile relative tolerance can waste rank or generate non-finite factors (screened kernels and spherical harmonics). A robust rule needs an absolute or global reference scale as well as floating-point safeguards.

\textbf{Failure signal.} More proxy points on the same set do not improve the answer (fractional/Riesz interactions), a two-sided kernel loses its row relations (Lippmann--Schwinger scattering), tiny matrices become subnormal (screened kernels and spherical harmonics), or a cheap representation costs too much to construct (the spherical-harmonic experiment's all-order precomputation).

\subsection{Choice 5: algebra and assembly}
\textbf{Question.} How do the local representations become an apply, transform or solve? Assembly is not interchangeable bookkeeping. It determines the data flow and the extra stability questions:
\begin{itemize}
 \item FMM translates and accumulates expansions without eliminating sources.
 \item Butterfly and FFT multiply a chain of sparse factors.
 \item Recursive skeletonization and FMM-LU sparsify, eliminate redundant unknowns, track Schur fill and finish with triangular/root solves.
\end{itemize}
Algebra does \emph{not} inherit stability from an accurate apply. Solving amplifies errors through conditioning and pivots. It also changes cost accounting: a smaller Schur root can coexist with larger storage and elimination work.

\textbf{Failure signal.} Pivots deteriorate even though $\|A_{\rm fast}-A\|$ is small (complex Helmholtz elimination), an approximate inverse fails at a caustic (the high-frequency-arc experiment), or a quantity of interest is accurate while the underlying density is not (the Oseen wake experiment).

\subsection{The algebra choice contains three assembly grammars}
\label{sec:grammars}
It is useful to make the three data flows explicit. They share block vocabulary but are not one universal algorithm.

\paragraph{G1: sum over levels (FMM and $\mathcal H^2$).}
Schematically,
\begin{equation}
 A\approx A_{\rm near}+\sum_{\ell=1}^{L}R_\ell^{\T}
 \left(\bigoplus_{\tau}T_{\ell,\tau}\right)S_\ell,
 \label{eq:g1}
\end{equation}
where $S_\ell$ forms source summaries, $T_{\ell,\tau}$ translates them across the interaction list, and $R_\ell$ distributes local information to targets. Upward merging, translation and downward inheritance are the algorithm.

\paragraph{G2: product of sparse factors (butterfly and FFT).}
The matrix becomes a chain such as
\begin{equation}
 A\approx U_LG_{L-1}\cdots G_hM_hH_h^*\cdots H_{L-1}^*V_L^*.
 \label{eq:g2}
\end{equation}
Application means applying the chain. Complementary low rank makes its local blocks small; nesting and radix recursion determine why the whole chain is cheap.

\paragraph{G3: eliminate (HBS, recursive skeletonization and FMM-LU).}
After row/column sparsification, redundant variables are eliminated. The basic Schur step is
\begin{equation}
 \begin{bmatrix}P&Q\\R&S\end{bmatrix}
 \begin{bmatrix}a\\b\end{bmatrix}=\begin{bmatrix}f\\g\end{bmatrix}
 \quad\Longrightarrow\quad
 (S-RP^{-1}Q)b=g-RP^{-1}f.
 \label{eq:schur}
\end{equation}
Repeating it bottom-up produces a direct factorization. Pivots, fill-in and modified blocks are now part of the representation, not an implementation afterthought.

\subsection{Three independent certificates}
A useful performance report maintains three separate budgets:
\begin{enumerate}
 \item \textbf{Approximation:} achieved error against a trusted reference at a stated tolerance and norm;
 \item \textbf{Stability:} safe pivots, conditioning and meaningful output errors for a solve;
 \item \textbf{Total cost:} construction, storage, application or solve work, and memory.
\end{enumerate}
A low rank is useful evidence for the first budget; it does not, by itself, certify the other two. A single accounting form is
\begin{equation}
 C\approx C_{\rm near}+C_{\rm build}
 +\sum_{\ell}n_\ell c_\ell C_{\rm transfer}(r_\ell)+C_{\rm root},
 \label{eq:costaccounting}
\end{equation}
where $n_\ell$ is the number of boxes at level $\ell$, $c_\ell$ the number of compressed partners per box, and $r_\ell$ their summary size. Every method needs this product, not just its $r_\ell$.

\section{From the same grammar to different algorithms}
\label{sec:algorithms}

\subsection{FMM: a tree of useful expansions}
The fast multipole method starts with a spatial tree, exact near interactions and separated far interactions. The expansion of a source box gives its outgoing summary; a translation sends it to another box; a local expansion evaluates its influence on target points. Its operation names have intuitive meanings:
\begin{itemize}
 \item sources to multipole (S2M): build a leaf's outgoing summary;
 \item multipole to multipole (M2M): re-express a child at its parent;
 \item multipole to local (M2L): translate an outgoing expansion into a local one;
 \item local to local (L2L): send inherited local information to children;
 \item local to target (L2T): evaluate potential where it is needed.
\end{itemize}
This is G1 from \eqref{eq:g1}. The logarithmic example above is a toy one-dimensional version of the same principle. The essential paper-specific contribution is \textbf{analytic representation and translation theory}: choices of centers, orders, transforms, error bounds and reuse. One tree and one admissibility rule do not derive those operators.

For a standard full-grid FMM same-level candidate list, the bound is $6^d-3^d$: 27 in two dimensions and 189 in three. With fixed expansion order in a non-oscillatory regime, the number of boxes decreases geometrically up the tree and the work can total $O(N)$. That is a statement about both the list and the transfer cost, not a consequence of the tree alone.

\subsection{Wideband FMM: change translation and representation with scale}
\label{sec:wideband}
Oscillatory kernels need more detail than the logarithmic example. Write a 3-D surface problem's top-level width in wavelengths as $W=kD/(2\pi)$ up to the report's convention; the essential scaling is $W\propto kD$. A large-box partial-wave order has leading form
\begin{equation}
 p=O\big(W+\log(1/\eps)\big),\qquad (p+1)^2\ \text{coefficients}.
\end{equation}
This leading large-box statement does not say that the separate low-frequency exponential expansions are accuracy-independent.

At high frequency, an ordinary expansion may be expensive to translate. A plane-wave or far-field-signature representation makes each direction-scale interaction almost \textbf{diagonal}: multiply each sampled direction independently. For a surface discretization with $N\propto(kD)^2$, the contrast is schematically
\begin{equation}
 \sum_\ell4^\ell(kw_\ell)^3=O(N^{3/2})
 \quad\text{versus}\quad
 \sum_\ell4^\ell(kw_\ell)^2=O(N\log N),
 \label{eq:wideband-cost}
\end{equation}
over the high-frequency levels; bounded-order finer levels add $O(N)$. Thus changing translation can matter more than changing rank.

At small scales, diagonal plane-wave forms can lose accuracy through cancellation and overflow before truncation. The wideband construction therefore combines low-frequency exponential/partial-wave machinery with high-frequency directional forms, conversion rules and filtering. The cited implementation's transition scales are $0.25$, $3.50$ and $12.0$ wavelengths at requested accuracies $10^{-3}$, $10^{-6}$ and $10^{-9}$, respectively \cite{cheng2006wideband}. They are measured transition choices, not one universal $\lambda/4$ rule. Its stated complexity is $O(N)$ when low-frequency work dominates and $O(N\log N)$ when high-frequency work dominates.

In our language:
\[
\text{FMM assembly}+\text{scale-appropriate expansions}
+\text{diagonal translations}+\text{a stable switch}.
\]
The major gain lies in representation and translation, not only in rank.

\subsection{Recursive skeletonization: use redundancy to eliminate}
\label{sec:rskel}
For unknowns $B$ inside a box, an ID chooses skeleton columns $S$ and redundant columns $R$ so that, against a far set $F$,
\begin{equation}
 A_{F,R}\approx A_{F,S}T.
 \label{eq:idrelation}
\end{equation}
Thus $S$ is a small set of interactions representing the rest against $F$. Nonsymmetric matrices need row as well as column control: the teaching driver stacks the two interaction directions rather than assuming symmetry.

Sparsification makes redundant unknowns couple only weakly to $F$. Block Gaussian elimination \eqref{eq:schur} then removes them, propagates skeletons upward and ends with a dense root solve. This is G3, not ``FMM run backwards'' in any literal sense \cite{ho2012hife,minden2017rss}. The new requirements are:
\begin{itemize}
 \item an elimination order that controls fill-in;
 \item stable pivots and conditioning estimates, not only small interpolation error;
 \item a list of blocks already changed by earlier Schur updates;
 \item singular and near-singular quadrature, which changes the original near entries.
\end{itemize}
In strong-admissibility FMM-LU, the modified-block or $Q$ set must include the appropriate non-kernel entries from the \emph{current} working matrix \cite{sushnikova2023fmmlu}. A proxy built only from the original kernel cannot see those Schur-generated entries unless the driver supplies them.

\clearpage
\begin{landscape}
\begin{table}[p]
\centering\footnotesize
\setlength{\tabcolsep}{3pt}
\caption{The complete algorithm map. ``G1/G2/G3'' refers to Section~\ref{sec:grammars}. The last column is the essential content that a generic block/tree interface does not supply.}
\label{tab:algorithm-map}
\begin{tabular}{@{}>{\raggedright\arraybackslash}p{1.55cm}>{\raggedright\arraybackslash}p{2.75cm}>{\raggedright\arraybackslash}p{2.65cm}>{\raggedright\arraybackslash}p{3.55cm}>{\raggedright\arraybackslash}p{3.45cm}>{\raggedright\arraybackslash}p{2.55cm}>{\raggedright\arraybackslash}p{5.85cm}@{}}
\toprule
Method & Partition / geometry & Interaction rule & Representation / compression & Sampling / construction & Assembly / grammar & Extra mathematics contributed by the algorithm \\
\midrule
FMM & One spatial tree shared by sources and targets & Exact near list; separated same-level far list & Outgoing/multipole and incoming/local analytic expansions & Kernel expansion formulas and quadrature; orders chosen by truncation bounds & \textbf{G1}: S2M, M2M, M2L, L2L, L2T & Addition/translation theorems, stable recurrences, expansion-order error bounds, interaction-list scaling and reuse. \\
Wideband FMM & FMM spatial tree with low- and high-frequency levels & Same near/far lists, but representation and regime switch depend on box size in wavelengths and $\eps$ & Low-frequency exponential/partial-wave forms; high-frequency plane-wave/far-field signatures; diagonal translators & Direction samples, interpolation/filtering and conversion between representation families & \textbf{G1} with scale-specific M2L and conversions & Stable diagonal translation, conversion/interpolation on the sphere, cancellation analysis and the precision-dependent switch. \\
HBS, RS, FMM-LU & One tree; nested skeletons per box & Weak off-diagonal or strong separated pairs; level-restricted strong lists; modified near blocks & ID skeletons $S,R$, interpolation $T$, Schur complements and LU factors & Exact true rows, proxy rows/volumes, sketches; current-$Q$ entries must be supplied in strong elimination & \textbf{G3}: sparsify, eliminate, triangle solves and dense root & Elimination order, fill-in/$Q$ bookkeeping, stable pivots, proxy representation theory and singular/near quadrature. \\
Butterfly & Two trees: row depth $\ell$ paired with column depth $L-\ell$ & Complementary phase-space tilings, not a distance predicate & Local low-rank tiles plus nested bases and small cross-level transfers & Full tile entries in our teacher; random products or sampled entries under separate hypotheses in fast constructors & \textbf{G2}: a product of $O(\log N)$ sparse factors & Phase regularity/complementary low rank, middle-level split, basis recursion, merge formulas and construction cost/incoherence theory. \\
FFT & Two index trees in mutually matched bit-reversed order & Complementary residue classes at every stage & Exact rank-one character tiles and diagonal twiddles & Exact index arithmetic; no statistical sampling & \textbf{G2}: Cooley--Tukey even/odd half-transforms and recombination & The residue identity \emph{and} radix recursion $T(N)=2T(N/2)+O(N)$, twiddle algebra and output permutation. \\
\bottomrule
\end{tabular}
\end{table}
\end{landscape}

\paragraph{Implementation boundary.} Table~\ref{tab:algorithm-map} is an algorithmic reading map, not a feature list for the included code. The Python package contains representative tree/admissibility/ID/proxy, recursive-solve and interpolative-butterfly drivers; it does not contain production FMM or wideband translation machinery. Rust contains smaller real-kernel examples, direct FFT/DFT-butterfly examples and a separate feasibility core. Section~\ref{sec:implementations} states exactly which library routine each shipped example reuses.

\takeaway{A shared data structure and measurement framework makes the algorithms easier to compare; it does not remove the expansions, translations, nested transfers, elimination theory or index recursion that make them work.}
